\begin{figure}[t]
\centering
\definecolor{pmsgc}{HTML}{2A78D6}
\definecolor{scigc}{HTML}{EB6834}
\definecolor{wfsgc}{HTML}{1BAF7A}
\begin{tikzpicture}[
  font=\normalsize, x=1cm, y=1cm, line cap=round, line join=round,
  frame/.style={draw=black!70, line width=0.8pt},
  hd/.style={anchor=south, inner sep=1pt},
  sub/.style={anchor=north, inner sep=1pt, text width=12.6cm, align=center}
]
%% One row per alternative: turbine rotor, shaft, generator (stator ring with the rotor symbol),
%% three-phase stator connection, back-to-back converter (ac/dc, dc link, dc/ac) and grid.
\newcommand{\chain}[1]{%
  \foreach \ang in {90,210,330} {
    \begin{scope}[shift={(1.1,#1)}, rotate=\ang-90]
      \path[frame, fill=black!15] (0.07,0) .. controls (0.18,0.22) and (0.13,0.52) .. (0,0.7)
                                   .. controls (-0.13,0.52) and (-0.18,0.22) .. (-0.07,0) -- cycle;
    \end{scope}
  }
  \fill[black!70] (1.1,#1) circle (0.09);
  \draw[black!55, line width=2.2pt] (1.25,#1) -- (3.7,#1);
  \filldraw[frame, fill=black!8] (4.3,#1) circle (0.62);
  \filldraw[frame, fill=white] (4.3,#1) circle (0.5);
  \draw[frame] (4.92,#1) -- (7.2,#1);
  \foreach \s in {-0.09,0,0.09} { \draw[frame] (6.0+\s,#1-0.12) -- (6.14+\s,#1+0.12); }
  \filldraw[frame, fill=black!5] (7.2,#1-0.4) rectangle (8.0,#1+0.4);
  \draw[frame] (7.2,#1-0.4) -- (8.0,#1+0.4);
  \node[inner sep=0] at (7.41,#1+0.2) {$\sim$};
  \node[inner sep=0] at (7.79,#1-0.2) {$=$};
  \draw[frame] (8.0,#1+0.1) -- (8.3,#1+0.1);
  \draw[frame] (8.0,#1-0.1) -- (8.3,#1-0.1);
  \filldraw[frame, fill=black!5] (8.3,#1-0.4) rectangle (9.1,#1+0.4);
  \draw[frame] (8.3,#1-0.4) -- (9.1,#1+0.4);
  \node[inner sep=0] at (8.51,#1+0.2) {$=$};
  \node[inner sep=0] at (8.89,#1-0.2) {$\sim$};
  \draw[frame] (9.1,#1) -- (11.3,#1);
  \draw[black!70, line width=2.4pt] (11.3,#1-0.5) -- (11.3,#1+0.5);
}
%% column headers
\node[hd] at (1.1,0.85) {turbine rotor};
\node[hd] at (4.3,0.85) {generator};
\node[hd] at (8.15,0.85) {converter};
\node[hd] at (11.3,0.85) {grid};
%% (a) SCIG: squirrel cage
\chain{0}
\draw[scigc, line width=1.1pt] (4.3,0) circle (0.42);
\draw[scigc, line width=1.1pt] (4.3,0) circle (0.26);
\foreach \a in {0,30,...,330} { \draw[scigc, line width=1.1pt] (4.3,0) ++(\a:0.26) -- ++(\a:0.16); }
\node[sub] at (6.2,-0.85) {(a) SCIG};
%% (b) WFSG: field winding on the rotor
\chain{-2.4}
\draw[black!50, line width=0.6pt] (4.3,-2.4) circle (0.42);
\draw[wfsgc, line width=1.1pt, decorate, decoration={coil, aspect=0.6, segment length=1.9mm, amplitude=1.9mm}] (3.9,-2.4) -- (4.7,-2.4);
\node[sub] at (6.2,-3.25) {(b) WFSG};
%% (c) PMSG: permanent magnets on the rotor surface
\chain{-4.8}
\draw[black!50, line width=0.6pt] (4.3,-4.8) circle (0.31);
\foreach \a in {0,120,240} { \draw[pmsgc, line width=2.8pt] (4.3,-4.8) ++(\a+7:0.4) arc[start angle=\a+7, end angle=\a+53, radius=0.4]; }
\foreach \a in {60,180,300} { \draw[black!40, line width=2.8pt] (4.3,-4.8) ++(\a+7:0.4) arc[start angle=\a+7, end angle=\a+53, radius=0.4]; }
\node[sub] at (6.2,-5.65) {(c) PMSG};
\end{tikzpicture}
\caption{The three alternatives as generator systems: turbine rotor, generator with its rotor excitation (squirrel cage, field winding, permanent magnets), back-to-back converter and grid}\label{fig:alternatives}
\end{figure}
